\documentclass[12pt]{article}
\usepackage[margin=1in]{geometry}
\usepackage{amsmath,amsthm,amssymb,scrextend}
\usepackage{array}
\setlength{\parindent}{0pt}

\begin{document}
\newcommand{\ds}{\displaystyle}



Let the horizontal plane containing the fixed point $P$ be the zero potential energy reference level (with the downward direction as the positive $y$-axis). The rope is uniform with a linear mass density of $\rho = \frac{2M}{2L} = \frac{M}{L}$.

\textbf{1. Initial State ($t=0$):}
\begin{itemize}
    \item Initial kinetic energy of the system: $E_{k0} = 0$
    \item Initial potential energy of the system (center of mass of the rope falls by $\frac{L}{2}$): $E_{p0} = -2Mg \left(\frac{L}{2}\right) = -MgL$
    \item \textbf{Total initial mechanical energy:} $E_0 = -MgL$
\end{itemize}

\textbf{2. Final State (particle has fallen a distance $x$):}
\begin{itemize}
    \item \textbf{Final kinetic energy:} The particle and the moving section of the rope have the same speed $v$.
    \[ E_{k} = \frac{1}{2}mv^2 + \frac{1}{2}\left[\frac{M}{L}\left(L - \frac{x}{2}\right)\right]v^2 = \frac{1}{2}\left(m + M - \frac{Mx}{2L}\right)v^2 \]
    \item \textbf{Final potential energy:} Composed of three parts: the particle, the stationary section of the rope (length $L + \frac{x}{2}$), and the moving section of the rope (length $L - \frac{x}{2}$).
    \[ E_{p} = -mgx - \underbrace{\frac{Mg}{L}\left(L + \frac{x}{2}\right)}_{\text{mass of stationary section}} \underbrace{\frac{1}{2}\left(L + \frac{x}{2}\right)}_{\text{COM of stationary section}} - \underbrace{\frac{Mg}{L}\left(L - \frac{x}{2}\right)}_{\text{mass of moving section}} \underbrace{\left(\frac{L}{2} + \frac{3x}{4}\right)}_{\text{COM of moving section}} \]
    Expanding and simplifying yields:
    \[ E_{p} = -mgx - MgL - Mgx + \frac{Mgx^2}{4L} \]
\end{itemize}

\textbf{3. Conservation of Mechanical Energy ($E_0 = E_k + E_p$):}
\[ -MgL = \frac{1}{2}\left(m + M - \frac{Mx}{2L}\right)v^2 - mgx - MgL - Mgx + \frac{Mgx^2}{4L} \]
\[ \frac{1}{2}\left(\frac{mL + ML - \frac{1}{2}Mx}{L}\right)v^2 = gx(m + M) - \frac{Mgx^2}{4L} \]
Extracting the common factor and rearranging proves the expression:
\[ v^2 = \frac{2gx(mL + ML - \frac{1}{4}Mx)}{mL + ML - \frac{1}{2}Mx} \]

\noindent\rule{\textwidth}{0.4pt}


Using the kinematic differential relationship $a = v\frac{dv}{dx} = \frac{1}{2}\frac{d(v^2)}{dx}$, we differentiate the above expression for $v^2$.

Let $v^2 = \frac{u}{w}$, where:
\begin{itemize}
    \item Numerator $u = 2gL(m+M)x - \frac{1}{2}Mgx^2 \implies u' = 2gL(m+M) - Mgx = 2gw$
    \item Denominator $w = L(m+M) - \frac{1}{2}Mx \implies w' = -\frac{1}{2}M$
\end{itemize}

According to the quotient rule:
\[ \frac{d(v^2)}{dx} = \frac{u'w - uw'}{w^2} = \frac{(2gw)w - u(-\frac{1}{2}M)}{w^2} = 2g + \frac{Mu}{2w^2} \]

Substituting into $a = \frac{1}{2}\frac{d(v^2)}{dx}$:
\[ a = g + \frac{Mu}{4w^2} \]

Substituting the original expressions for $u$ and $w$ back and simplifying proves the expression:
\[ a = g + \frac{Mgx(mL + ML - \frac{1}{4}Mx)}{2\left(mL + ML - \frac{1}{2}Mx\right)^2} \]

\noindent\rule{\textwidth}{0.4pt}

Observe the expression for acceleration:
\[ a = g + \frac{Mgx(mL + ML - \frac{1}{4}Mx)}{2\left(mL + ML - \frac{1}{2}Mx\right)^2} \]

At any instant after the particle is released ($x > 0$):
\begin{enumerate}
    \item \textbf{Denominator:} $2\left(mL + ML - \frac{1}{2}Mx\right)^2$ is a perfect square and is always strictly positive ($> 0$).
    \item \textbf{Numerator:} The problem states that $x < 2L$, therefore $\frac{1}{4}Mx < \frac{1}{4}M(2L) = \frac{1}{2}ML$.
    From this, we can deduce the term inside the parenthesis:
    \[ (mL + ML - \frac{1}{4}Mx) > \left(mL + ML - \frac{1}{2}ML\right) = mL + \frac{1}{2}ML > 0 \]
\end{enumerate}

Since $M, g, x$ are all positive numbers, and the term inside the parenthesis is also positive, the entire fractional term is strictly greater than $0$. Therefore:
\[ a = g + (\text{positive number}) > g \]
The original proposition is proven.

\bigskip
\hrule
\begin{center}
    This solution was submitted by \textbf{Roy}\\[4pt]
    Published and refined by \textbf{Jon}\\[8pt]
    \textit{Last updated: July 25, 2026}
\end{center}

\end{document}